\chapter{Functional evaluation}
\label{chap:functional-evaluation}

\begin{quote}
Spec dir:
\passthrough{\lstinline!specs/2026-05-02-phase-4-functional-evaluation/!}.
Version: \passthrough{\lstinline!0.4.0!}.
\end{quote}

\subsection{What this gives you}\label{what-this-gives-you}

A primitive that evaluates a \textbf{functional} of a scalar field ---
an integral of an arbitrary user-supplied integrand
\(f(\psi, \nabla\psi, \nabla^2\psi)\) over the periodic Cartesian
domain:

\[
F[\psi] \;=\; \int_{\Omega} f\bigl(\psi(\mathbf{x}),\, \nabla\psi(\mathbf{x}),\, \nabla^2\psi(\mathbf{x})\bigr)\, dV.
\]

\passthrough{\lstinline!evaluate!} builds on top of Phase 1's
\passthrough{\lstinline!laplacian!}, Phase 2's
\passthrough{\lstinline!gradient!}, and Phase 3's
\passthrough{\lstinline!integrate!}: it materializes whichever
derivatives your \passthrough{\lstinline!f!} actually needs, evaluates
\passthrough{\lstinline!f!} pointwise into a temporary
\passthrough{\lstinline!Field<double>!}, and reduces the result with the
same pairwise / Kahan tag dispatch from Phase 3.

\subsection{API}\label{api}

Header:
\href{../../../include/gridcalc/func/Functional.hpp}{\passthrough{\lstinline!include/gridcalc/func/Functional.hpp!}}.

\begin{lstlisting}[language={C++}]
namespace gridcalc::func {

template <class F, class Tag = Pairwise>
double evaluate(const core::Field<double>& psi, F&& f, Tag tag = {});

}  // namespace gridcalc::func
\end{lstlisting}

Your callable \passthrough{\lstinline!f!} may take \textbf{any one of
three signatures}, picked at compile time by
\passthrough{\lstinline!std::is\_invocable\_r\_v!}:

{\def\LTcaptype{none} % do not increment counter
\begin{longtable}[]{@{}
  >{\raggedright\arraybackslash}p{(\linewidth - 2\tabcolsep) * \real{0.8361}}
  >{\raggedright\arraybackslash}p{(\linewidth - 2\tabcolsep) * \real{0.1639}}@{}}
\toprule\noalign{}
\begin{minipage}[b]{\linewidth}\raggedright
Signature
\end{minipage} & \begin{minipage}[b]{\linewidth}\raggedright
Computes
\end{minipage} \\
\midrule\noalign{}
\endhead
\bottomrule\noalign{}
\endlastfoot
\passthrough{\lstinline!f(double psi)!} & only
\passthrough{\lstinline!ψ!} --- gradient and Laplacian are
\textbf{skipped} \\
\passthrough{\lstinline!f(double psi, const Vec3d\& grad)!} &
\passthrough{\lstinline!ψ!} and \passthrough{\lstinline!∇ψ!} ---
Laplacian is \textbf{skipped} \\
\passthrough{\lstinline!f(double psi, const Vec3d\& grad, double lap)!}
& full triplet \passthrough{\lstinline!ψ!},
\passthrough{\lstinline!∇ψ!}, \passthrough{\lstinline!∇²ψ!} \\
\end{longtable}
}

Return type must be convertible to \passthrough{\lstinline!double!}. The
reducer tag (\passthrough{\lstinline!Pairwise\{\}!} or
\passthrough{\lstinline!Kahan\{\}!}) is forwarded to
\passthrough{\lstinline!func::integrate!}.

\subsection{Worked example: Ginzburg--Landau free
energy}\label{worked-example-ginzburglandau-free-energy}

The Ginzburg--Landau free-energy density is

\[
f_{\mathrm{GL}}(\psi, \nabla\psi) \;=\; \tfrac{1}{2} |\nabla\psi|^2 \;+\; W (\psi^2 - 1)^2,
\]

with \(W\) a positive constant. Note that \(f_{\mathrm{GL}}\) does
\textbf{not} depend on \(\nabla^2\psi\), so writing it with the 2-arg
signature lets \passthrough{\lstinline!evaluate!} skip the Laplacian
computation.

\subsubsection{Setup}\label{setup}

\begin{lstlisting}[language={C++}]
#include <cmath>
#include <gridcalc/core/EigenAliases.hpp>
#include <gridcalc/core/Field.hpp>
#include <gridcalc/core/Grid.hpp>
#include <gridcalc/func/Functional.hpp>

using gridcalc::core::Field;
using gridcalc::core::Grid;
using gridcalc::core::Vec3d;
using gridcalc::func::evaluate;

constexpr double kPi = 3.14159265358979323846;

const int N = 64;
const double L = 2.0 * kPi;
const double h = L / N;
const double W = 1.0;

Grid grid(N, N, N, Vec3d(h, h, h));
Field<double> psi(grid, [](double x, double y, double z) {
  return std::sin(x) * std::sin(y) * std::sin(z);
});
\end{lstlisting}

\subsubsection{The integrand}\label{the-integrand}

\begin{lstlisting}[language={C++}]
auto gl_density = [W](double p, const Vec3d& g) {
  const double dpsi2 = p * p - 1.0;
  return 0.5 * g.squaredNorm() + W * dpsi2 * dpsi2;
};
\end{lstlisting}

\subsubsection{Evaluate}\label{evaluate}

\begin{lstlisting}[language={C++}]
const double F = evaluate(psi, gl_density);
\end{lstlisting}

For this \passthrough{\lstinline!ψ!} with
\passthrough{\lstinline!W = 1!}, the analytical reference is
\(F = (3/2)\pi^3 + (411/64)\pi^3 = (507/64)\pi^3 \approx 245.62\). At
\passthrough{\lstinline!N = 64!}, \passthrough{\lstinline!evaluate!}
matches this to better than 0.5\% relative error; the discrete error is
dominated by the squared-gradient term and converges at order 2 in \(h\)
(verified by
\href{../../../test/functional_test.cpp}{\passthrough{\lstinline!test/functional\_test.cpp!}}'s
\passthrough{\lstinline!GinzburgLandauSecondOrderConvergence!} test).

For the why-this-formula-converges-at-order-2 background --- derivation,
error analysis, references --- see
\href{../../developer-note/notes/phase-4-functional-evaluation.md}{\passthrough{\lstinline!docs/developer-note/notes/phase-4-functional-evaluation.md!}}.

\subsection{\texorpdfstring{What this does \emph{not} do
(yet)}{What this does not do (yet)}}\label{what-this-does-not-do-yet}

\begin{itemize}
\tightlist
\item
  \textbf{No vector-field functionals.}
  \passthrough{\lstinline!evaluate(Field<Vec3d>, …)!} arrives in Phase
  11+ when the higher-rank functionals first need it.
\item
  \textbf{No higher-order derivatives in \passthrough{\lstinline!f!}.}
  Phase 6 adds biharmonic and mixed partials; Phase 11 adds
  up-to-4th-order gradients.
\item
  \textbf{No variational derivative \passthrough{\lstinline!δF/δψ!}.}
  Phase 12 (CH demo) adds closed-form variational derivatives where they
  are needed.
\item
  \textbf{No fused-loop optimization.}
  \passthrough{\lstinline!evaluate!} materializes intermediate Fields
  for \passthrough{\lstinline!∇ψ!} and/or \passthrough{\lstinline!∇²ψ!}.
  Phase 11+ replaces this with expression-template fusion.
\end{itemize}

\subsection{Quick reference: which arity to
pick}\label{quick-reference-which-arity-to-pick}

\begin{itemize}
\tightlist
\item
  Integrand depends only on \passthrough{\lstinline!ψ!} (e.g.,
  \passthrough{\lstinline!∫ ψ² dV!},
  \passthrough{\lstinline!∫ W(ψ²-1)² dV!}): use
  \passthrough{\lstinline!f(double)!}. Saves both stencil computations.
\item
  Integrand uses \passthrough{\lstinline!∇ψ!} (e.g., the gradient term
  of any GL-style energy): use
  \passthrough{\lstinline!f(double, const Vec3d\&)!}. Saves the
  Laplacian.
\item
  Integrand uses \passthrough{\lstinline!∇²ψ!} (Cahn--Hilliard kernel,
  surface-tension derivatives): use
  \passthrough{\lstinline!f(double, const Vec3d\&, double)!}.
\end{itemize}
